% =====================================================================
%  Sink Detection
%
%  Self-contained. Every construction used is defined here; no companion
%  paper is required. External citations are to the mass spectrometry,
%  graph theory, and algorithms literature only.
% =====================================================================
\documentclass[11pt,a4paper]{article}

\usepackage[utf8]{inputenc}
\usepackage[T1]{fontenc}
\usepackage{amsmath,amssymb,amsfonts,amsthm}
\usepackage{mathtools}
\usepackage[margin=1in]{geometry}
\usepackage{graphicx}
\usepackage{booktabs}
\usepackage{array}
\usepackage{enumitem}
\usepackage{lmodern}
\usepackage[protrusion=true,expansion=false]{microtype}
\usepackage[numbers,sort&compress]{natbib}
\usepackage[colorlinks=true,linkcolor=blue,citecolor=blue,urlcolor=blue]{hyperref}
\usepackage[capitalise,nameinlink]{cleveref}

% ---------------------------------------------------------------------
%  Theorem environments
% ---------------------------------------------------------------------
\newtheorem{theorem}{Theorem}[section]
\newtheorem{lemma}[theorem]{Lemma}
\newtheorem{corollary}[theorem]{Corollary}
\newtheorem{proposition}[theorem]{Proposition}
\theoremstyle{definition}
\newtheorem{definition}[theorem]{Definition}
\newtheorem{axiom}[theorem]{Axiom}
\newtheorem{assumption}[theorem]{Assumption}
\newtheorem{construction}[theorem]{Construction}
\newtheorem{example}[theorem]{Example}
\theoremstyle{remark}
\newtheorem{remark}[theorem]{Remark}

% ---------------------------------------------------------------------
%  Notation
% ---------------------------------------------------------------------
\newcommand{\Gph}{\mathcal{G}}
\newcommand{\wt}{w}
\newcommand{\med}{\mathfrak{m}}
\newcommand{\flo}{\beta}
\newcommand{\cut}{\operatorname{cut}}
\newcommand{\sep}{\sigma}
\newcommand{\dep}{\delta}
\newcommand{\spr}{\rho}                       % spread
\newcommand{\Snk}{\mathsf{Sink}}
\newcommand{\att}{a}                          % attachment
\DeclareMathOperator*{\argmin}{arg\,min}
\DeclareMathOperator*{\argmax}{arg\,max}

\title{\textbf{Sink Detection:\\[2pt]
Finding the Vertices That Absorb Every Question}}

\author{Kundai Farai Sachikonye\\
\texttt{kundai.sachikonye@wzw.tum.de}}

\date{\today}

\begin{document}
\maketitle

% =====================================================================
\begin{abstract}
\noindent
In a weighted contact graph built from analytical data, a small number of
vertices are adjacent to almost everything. A ubiquitous contaminant ion,
a solvent cluster, a badly-behaved column artefact, or a term that occurs
in every record of a reference library all produce the same structure:
a vertex whose neighbourhood is a large fraction of the graph. We call
such a vertex a \emph{sink}. Sinks are destructive in a way that is not
obvious, and the purpose of this paper is to make the damage precise and
give a cheap test. We show that a single sink of sufficient attachment
collapses the separation structure of the entire graph, driving every
query's answer to within a bounded band of every other's
(\Cref{thm:collapse}), that this happens without any query returning an
error or an implausible value (\Cref{cor:silent}), and that the effect is
not diluted by graph size but amplified by it (\Cref{prop:amplify}). We
then give a detection criterion. Sinks are \emph{not} identified by
degree, and we prove this with a counterexample family in which a
high-degree vertex is harmless and a moderate-degree vertex is fatal
(\Cref{thm:degree-fails}). The correct statistic is \emph{spread} --- the
fraction of the graph a vertex separates from --- and we prove that
thresholding spread is sound (\Cref{thm:spread-sound}), computable in one
pass (\Cref{prop:spread-cost}), and that the threshold is forced rather
than tuned, being pinned by the separation floor
(\Cref{thm:threshold}). Finally we prove that excision is the correct
remedy and reweighting is not (\Cref{thm:excision}), and give the
diagnostic protocol.
\end{abstract}

% =====================================================================
\section{Introduction}
\label{sec:intro}

Build a weighted graph in which vertices are observations from an
analytical experiment and edges record measured contact --- mass
proximity, co-elution, shared fragments, shared library annotations.
Almost every such graph, built from almost any real dataset, contains a
few vertices adjacent to an enormous fraction of the rest.

The sources are mundane and various. A plasticiser ion appears in every
scan. A solvent cluster series lands within tolerance of most low-mass
features. A retention-time artefact co-elutes with everything by
construction. In a reference library, an annotation term such as
\emph{organic compound} is attached to every record. In each case a
single vertex acquires a neighbourhood comprising a large fraction of the
graph.

The natural reaction is that this is a nuisance, handled by filtering
obvious contaminants. We argue it is worse than a nuisance and that the
handling is not obvious. Our claim is that one such vertex is sufficient
to make every separation computation in the graph return an answer that
is nearly independent of what was asked, and that this happens
\emph{silently}: no computation fails, no value is out of range, and the
returned numbers remain plausible and even reproducible. A pipeline
contaminated by a sink produces a stable, reproducible, meaningless
output, which is worse than an unstable one because it survives quality
control.

This paper makes that claim a theorem, shows that the obvious detector
(degree) does not work, gives one that does, and proves the remedy.

\subsection{What is proved}

\Cref{sec:setting} constructs the contact graph and the separation
functional from scratch and establishes the floor
(\Cref{thm:floor}). \Cref{sec:damage} defines sinks and proves the
collapse theorem (\Cref{thm:collapse}), its silence
(\Cref{cor:silent}), and its amplification with size
(\Cref{prop:amplify}). \Cref{sec:degree} proves that degree is not a
detector (\Cref{thm:degree-fails}). \Cref{sec:spread} defines spread,
proves soundness (\Cref{thm:spread-sound}), cost
(\Cref{prop:spread-cost}), and that the threshold is forced
(\Cref{thm:threshold}). \Cref{sec:remedy} proves excision is correct and
reweighting is not (\Cref{thm:excision},
\Cref{prop:reweight-fails}). \Cref{sec:protocol} gives the protocol.

% =====================================================================
\section{Setting}
\label{sec:setting}

\begin{definition}[Contact graph]\label{def:contact}
A \emph{contact graph} is a finite weighted graph $\Gph = (V,E,\wt)$ with
$\wt : E \to \mathbb{R}_{>0}$ and a distinguished vertex $\med \in V$,
the \emph{medium}, adjacent to every other vertex. We write
$U = V \setminus \{\med\}$ for the \emph{items}, $n = |U|$, and for
$S \subseteq V$
\[
\cut(S) \;=\; \sum_{\{u,v\} \in E,\ |S \cap \{u,v\}| = 1} \wt(\{u,v\}).
\]
\end{definition}

\begin{remark}[Why the medium is definitional]\label{rem:medium}
The medium is not an item and does not correspond to a measurement. It
represents ``everything else'', and $\wt(\{u,\med\})$ is the cost of
asserting that $u$ is a distinct thing rather than an aspect of the
background. Its adjacency to every vertex is part of
\Cref{def:contact}, not a property of the data: if items were
individuated only against nominated rivals, the nomination of each rival
would require a prior individuation and the argument would regress. The
medium terminates the regress by being adjacent by definition. It cannot
be deleted, and \Cref{thm:floor} fails without it.
\end{remark}

\begin{definition}[Separation and depth]\label{def:sep}
For $v \in U$, an \emph{admissible set} is $S \subseteq V$ with $v \in S$
and $\med \notin S$. Put
\[
S^{\ast}(v) = \argmin\{\cut(S) : S \text{ admissible}\},
\qquad
\sep(v) = \cut(S^{\ast}(v)),
\qquad
\dep(v) = |S^{\ast}(v)| ,
\]
taking the inclusion-minimal minimiser when several attain the minimum.
\end{definition}

\begin{proposition}[Computability]\label{prop:computable}
$\sep(v)$ and $S^{\ast}(v)$ are computed by a single maximum-flow
computation with source $v$ and sink $\med$, in $O(|V||E|)$ time.
\end{proposition}

\begin{proof}
Admissible sets are exactly the $v$--$\med$ cuts, so $\sep(v)$ is the
minimum $v$--$\med$ cut value, which equals the maximum flow value
\citep{fordfulkerson1956,menger1927}. Push--relabel computes a maximum
flow in $O(|V||E|)$ \citep{goldberg1988new}, and the inclusion-minimal
minimum cut is read from the residual graph as the set reachable from
$v$.
\end{proof}

\begin{axiom}[Non-completability]\label{ax:noncomplete}
For every item set $U$ arising from an experiment there is a conceivable
extension $U' \supsetneq U$, and no procedure internal to $U$ certifies
that no such extension exists.
\end{axiom}

\begin{theorem}[Separation floor]\label{thm:floor}
Under \Cref{ax:noncomplete} there is $\flo > 0$ with $\sep(v) \geq \flo$
for every $v \in U$.
\end{theorem}

\begin{proof}
Every admissible $S$ for $v$ contains $v$ and excludes $\med$, so the
edge $\{v,\med\}$, present by \Cref{def:contact}, crosses $S$. Hence
$\sep(v) \geq \wt(\{v,\med\}) \geq \min_{u \in U}\wt(\{u,\med\}) =:
\flo_0 > 0$, positive because $\wt$ is $\mathbb{R}_{>0}$-valued and $U$
is finite. Setting $\flo_0 = 0$ would assert that some item is separable
from everything at no cost, that is, individuated by inspection of itself
alone; by \Cref{ax:noncomplete} no internal procedure certifies that no
further item shares its description, so that assignment is unlicensed.
Take $\flo = \flo_0$.
\end{proof}

% =====================================================================
\section{What a Sink Does}
\label{sec:damage}

\begin{definition}[Attachment and sink]\label{def:sink}
Let $z \in U$. Its \emph{attachment} is
\[
\att(z) \;=\; \min_{u \in U \setminus \{z\}} \wt(\{u,z\})
\quad\text{over the edges present},
\]
and its \emph{neighbourhood fraction} is
$|N(z) \cap U| / (n-1)$, where $N(z)$ is the set of vertices adjacent to
$z$. We call $z$ a \emph{sink at level $(\lambda,\theta)$} if
$\att(z) \geq \lambda$ and its neighbourhood fraction is at least
$\theta$.
\end{definition}

\begin{theorem}[Collapse]\label{thm:collapse}
Let $z$ be adjacent to every item, with $\wt(\{u,z\}) \geq \lambda$ for
all $u \in U \setminus \{z\}$. Then for every $v \in U \setminus \{z\}$,
\[
\flo \;\leq\; \sep(v) \;\leq\; \wt(\{v,\med\}) + \wt(\{v,z\}) + D(v),
\]
where $D(v)$ is the total weight of $v$'s non-medium, non-$z$ edges; and
moreover, for every $v$ whose separating set does not contain $z$,
\[
\sep(v) \;\geq\; \lambda \cdot \dep(v).
\]
Consequently, if the separating sets have depth at least $d$ and
$\lambda d$ exceeds the spread of the medium weights, all separations lie
in a band of width determined by $\lambda$ rather than by any structure
peculiar to $v$.
\end{theorem}

\begin{proof}
The lower bound $\sep(v)\geq\flo$ is \Cref{thm:floor}. For the upper
bound, take the admissible set $S = \{v\}$: its cut consists of the
medium edge, the edge to $z$, and $v$'s remaining edges, so
$\sep(v) \leq \cut(\{v\}) = \wt(\{v,\med\}) + \wt(\{v,z\}) + D(v)$.

For the second display, let $S$ be an admissible minimiser with
$z \notin S$. Every $u \in S$ is adjacent to $z$ by hypothesis, and
$z \notin S$, so each edge $\{u,z\}$ crosses $S$. These edges are
distinct for distinct $u$, and each has weight at least $\lambda$, so
\[
\sep(v) \;=\; \cut(S) \;\geq\; \sum_{u \in S \cap U} \wt(\{u,z\})
\;\geq\; \lambda \cdot |S| \;=\; \lambda\,\dep(v),
\]
where $|S| = \dep(v)$ since $S$ is the minimiser.

For the consequence: if $\dep(v) \geq d$ for all $v$, then
$\sep(v) \in [\lambda d,\ \lambda\,\dep(v) + \wt(\{v,\med\}) + D(v)]$,
and when $\lambda d$ dominates the medium and residual terms the interval
endpoints are governed by $\lambda$ and $\dep(v)$, both of which are
properties of $z$'s attachment rather than of $v$'s own contacts. Two
items with entirely different contact structures then receive separations
differing by at most the variation in the subordinate terms.
\end{proof}

\begin{corollary}[The collapse is silent]\label{cor:silent}
Under the hypotheses of \Cref{thm:collapse}, every $\sep(v)$ is a
positive real number, no computation is undefined, the values are stable
under recomputation, and the values are reproducible across runs of the
same acquisition. No internal test of the returned numbers distinguishes
this situation from an uncontaminated graph.
\end{corollary}

\begin{proof}
Positivity is \Cref{thm:floor}. Definedness holds because a minimum over
a non-empty finite collection of admissible sets always exists ---
$\{v\}$ is always admissible. Stability and reproducibility hold because
$\sep$ is a deterministic function of $(\Gph,\wt)$ and the graph is fixed
by the data. For the last claim: the range of returned values in the
contaminated graph is an interval of positive reals, and an interval of
positive reals is also the range in an uncontaminated graph, so no
predicate on the multiset of returned values alone separates the cases.
\end{proof}

\begin{remark}[Why silence is the whole problem]\label{rem:silence}
\Cref{cor:silent} is why sinks warrant a paper. A computation that failed
loudly would be caught by any quality control worth the name. A
computation that returns stable, reproducible, positive, plausible
numbers that happen not to depend on the input passes every check
designed to catch instability, and passes it repeatedly, and passes it
across laboratories. Reproducibility is not evidence against
contamination here; it is a prediction of it.
\end{remark}

\begin{proposition}[Amplification with size]\label{prop:amplify}
Under the hypotheses of \Cref{thm:collapse}, the contribution of $z$ to
$\cut(S)$ for an admissible $S$ omitting $z$ grows linearly in $|S|$,
while the contribution of $v$'s own non-$z$ contacts is bounded by $D(v)$
independently of $n$. Hence the fraction of $\sep(v)$ attributable to $z$
tends to $1$ as $\dep(v)$ grows.
\end{proposition}

\begin{proof}
From the proof of \Cref{thm:collapse}, the $z$-edges crossing $S$
contribute at least $\lambda|S|$. The non-$z$ contributions of $v$ are at
most $D(v) + \wt(\{v,\med\})$ plus the corresponding bounded terms of the
other members of $S$; if each item's non-$z$ contact weight is bounded by
a constant $C$ independent of $n$, the non-$z$ total is at most $C|S|$
while the $z$ total is at least $\lambda|S|$, so the $z$-fraction is at
least $\lambda/(\lambda+C)$ for every $|S|$, and taking $\lambda$ large
or the analysis to deeper $S$ drives it to $1$.
\end{proof}

\begin{remark}[The wrong intuition]\label{rem:dilution}
One might expect a single bad vertex to matter less as the graph grows,
being one of $n$. \Cref{prop:amplify} says the opposite: because the sink
is adjacent to \emph{everything}, it contributes an edge for every member
of every separating set, so its total contribution scales with the size
of what is being separated while a normal vertex's does not. Collecting
more data makes a sink worse, not better.
\end{remark}

% =====================================================================
\section{Degree Is Not a Detector}
\label{sec:degree}

The obvious test is to flag high-degree vertices. It does not work, and
the failure is in both directions.

\begin{theorem}[Degree fails]\label{thm:degree-fails}
There exist contact graphs containing a vertex of degree $\Theta(n)$ that
is harmless --- its deletion changes no $\sep(v)$ by more than an
arbitrarily small $\epsilon$ --- and a vertex of degree
$\Theta(\sqrt{n})$ whose deletion changes $\sep(v)$ by a factor bounded
away from $1$ for a constant fraction of $v$. Hence no threshold on
degree is both sound and complete for sinkhood.
\end{theorem}

\begin{proof}
\emph{Harmless high degree.} Take $U$ of size $n$, medium edges of weight
$1$, and a vertex $h$ joined to all of $U \setminus \{h\}$ with weight
$\epsilon/n$. Then $\deg(h) = n-1 = \Theta(n)$. For any $v$ and any
admissible $S$, the $h$-edges crossing $S$ contribute at most
$(\epsilon/n)\cdot n = \epsilon$. Deleting $h$ therefore changes every
cut value by at most $\epsilon$, hence changes $\sep(v)$ by at most
$\epsilon$.

\emph{Fatal moderate degree.} Take $U$ partitioned into
$\sqrt{n}$ blocks of size $\sqrt{n}$, with medium edges of weight $1$,
no contact edges within or between blocks, and a vertex $z$ joined with
weight $\lambda \gg 1$ to exactly one designated block $B$, so
$\deg(z) = \sqrt{n} = \Theta(\sqrt{n})$. For $v \in B$: with $z$ present,
$\cut(\{v\}) = 1 + \lambda$, and any larger admissible $S \subseteq B$
pays $\lambda$ per member, so $\sep(v) = 1+\lambda$. With $z$ deleted,
$\sep(v) = 1$. The ratio is $1+\lambda$, bounded away from $1$, and $B$
has $\sqrt{n}$ members --- a constant fraction of $n$ if the construction
is repeated across a constant fraction of blocks, which changes only
$\deg(z)$ by a constant factor.

Since the harmless vertex has strictly larger degree than the fatal one
for $n$ sufficiently large, any threshold flagging the fatal vertex also
flags the harmless one, and any threshold sparing the harmless one also
spares the fatal one.
\end{proof}

\begin{remark}[What degree misses]\label{rem:degree}
Degree counts edges and ignores both their weights and their effect on
cuts. \Cref{thm:degree-fails} shows both omissions matter: the first
construction has many edges of negligible weight, the second few edges of
decisive weight. The statistic that matters is not how many vertices $z$
touches but how much of the graph's separation structure passes through
it.
\end{remark}

% =====================================================================
\section{Spread}
\label{sec:spread}

\begin{definition}[Spread]\label{def:spread}
For $z \in U$, the \emph{spread} of $z$ is
\[
\spr(z) \;=\; \frac{1}{n-1}\;
\bigl|\{\, v \in U \setminus \{z\} \;:\;
z \in \partial S^{\ast}(v) \,\}\bigr| ,
\]
where $\partial S$ denotes the set of vertices incident to an edge
crossing $S$. That is, $\spr(z)$ is the fraction of items whose
separating cut $z$ actually crosses.
\end{definition}

\begin{remark}[Spread is not degree]\label{rem:spread-vs-degree}
Spread counts the separations $z$ participates in, not the vertices $z$
touches. In the harmless construction of \Cref{thm:degree-fails}, $h$ has
degree $n-1$; but for $\epsilon$ small the minimisers $S^{\ast}(v)$ are
unchanged by $h$'s presence, and while $h$ is incident to crossing edges
its \emph{weighted} participation is $\epsilon$ --- which motivates the
weighted variant used in \Cref{thm:spread-sound}.
\end{remark}

\begin{definition}[Weighted spread]\label{def:wspread}
The \emph{weighted spread} of $z$ is
\[
\spr_{\wt}(z) \;=\; \frac{1}{n-1}
\sum_{v \in U \setminus \{z\}}
\frac{\sum_{e \ni z,\ e \text{ crosses } S^{\ast}(v)} \wt(e)}{\sep(v)} ,
\]
the mean fraction of each separation cost contributed by edges at $z$.
By construction $0 \leq \spr_{\wt}(z) \leq 1$.
\end{definition}

\begin{theorem}[Spread is sound]\label{thm:spread-sound}
Let $z \in U$ and suppose $\spr_{\wt}(z) \leq \tau$. Then for every
$v \in U \setminus \{z\}$, deleting $z$ changes $\sep(v)$ by at most
$\tau' \sep(v)$ where $\tau'$ is the per-$v$ contribution fraction, and
in the mean over $v$ the change is at most $\tau$ in relative terms.
Conversely, if $z$ is a sink at level $(\lambda,1)$ in the sense of
\Cref{def:sink} and every $\dep(v) \geq d$, then
\[
\spr_{\wt}(z) \;\geq\; \frac{\lambda d}{\lambda d + C d}
\;=\; \frac{\lambda}{\lambda + C},
\]
with $C$ the per-item bound on non-$z$ contact weight from
\Cref{prop:amplify}. Hence sinks have large weighted spread and vertices
of small weighted spread are not sinks.
\end{theorem}

\begin{proof}
For the first claim, fix $v$ and let $S = S^{\ast}(v)$. Deleting $z$
removes exactly the $z$-incident edges crossing $S$, whose total weight
is by definition $\tau'_v \sep(v)$ with $\tau'_v$ the summand in
\Cref{def:wspread}. The cut value of $S$ in the reduced graph is
therefore $(1-\tau'_v)\sep(v)$, and since deletion cannot increase any
cut value, the new separation is at most this and at least
$\sep(v) - \tau'_v\sep(v)$ --- the latter because any admissible set in
the reduced graph has cut value at least its value in $\Gph$ minus the
$z$-edges it crossed, which is at most the maximum $z$-contribution.
Averaging over $v$ gives the mean statement with $\tau$.

For the converse, by \Cref{thm:collapse} the $z$-edges crossing
$S^{\ast}(v)$ contribute at least $\lambda\dep(v) \geq \lambda d$, while
$\sep(v) \leq \lambda\dep(v) + C\dep(v)$ by the same argument applied to
the bounded non-$z$ contacts. The ratio is therefore at least
$\lambda d/(\lambda d + Cd)$ for every $v$, and the mean of quantities
each at least this bound is at least this bound.
\end{proof}

\begin{proposition}[Cost]\label{prop:spread-cost}
Given the $n$ minimisers $S^{\ast}(v)$ --- one max-flow each, $O(n|V||E|)$
in total --- all $n$ weighted spreads are computed in one further pass in
$O(n|E|)$ time and $O(|V|)$ additional space, by accumulating, for each
$v$, the weight of each crossing edge into the accumulators of its two
endpoints.
\end{proposition}

\begin{proof}
For each $v$, iterate over the edges crossing $S^{\ast}(v)$ --- at most
$|E|$ of them --- and for each such edge $e = \{x,y\}$ add
$\wt(e)/\sep(v)$ to the accumulators of $x$ and of $y$. After all $v$,
divide each accumulator by $n-1$. Each edge is examined once per $v$,
giving $O(n|E|)$, and one accumulator per vertex gives $O(|V|)$ space.
\end{proof}

\begin{theorem}[The threshold is forced]\label{thm:threshold}
Let $\flo$ be the floor of \Cref{thm:floor} and let $W = \max_{v}\sep(v)$.
If $\spr_{\wt}(z) > 1 - \flo/W$ then for some $v$ the deletion of $z$
would leave a separation below the floor, contradicting
\Cref{thm:floor} for the reduced graph unless the reduced graph's own
floor is smaller. Consequently the admissible threshold on weighted
spread is not a free parameter: it is bounded above by $1 - \flo/W$,
a quantity determined by the graph.
\end{theorem}

\begin{proof}
Suppose $\spr_{\wt}(z) > 1 - \flo/W$. Then for at least one $v$ the
per-$v$ fraction satisfies $\tau'_v > 1 - \flo/W$, so the $z$-contribution
to $\sep(v)$ exceeds $(1-\flo/W)\sep(v)$, leaving at most
$(\flo/W)\sep(v) \leq (\flo/W)W = \flo$ of $\sep(v)$ attributable to
anything else. By the first part of \Cref{thm:spread-sound}, deleting
$z$ leaves $S^{\ast}(v)$ with cut value at most $\flo$, so the reduced
graph's separation for $v$ is at most $\flo$. Since \Cref{thm:floor}
holds in the reduced graph with its own floor $\flo'$, we get
$\flo' \leq \flo$, meaning $z$'s removal has lowered the floor --- $z$
was supplying it. A vertex supplying the floor for another item is
precisely a vertex whose contact is doing the individuating work that the
medium is supposed to do, which is the defining pathology. Hence values
of the threshold above $1-\flo/W$ admit such vertices and the threshold
must lie below it.
\end{proof}

\begin{remark}[Why forcing matters]\label{rem:forced}
A detector with a tunable threshold is a detector whose output can be
argued with. \Cref{thm:threshold} pins the admissible range using $\flo$
and $W$, both computed from the graph in the same pass that computes the
spreads. The practitioner does not choose the cutoff; the graph reports
it.
\end{remark}

% =====================================================================
\section{The Remedy}
\label{sec:remedy}

Two remedies suggest themselves: delete the sink, or down-weight its
edges. We prove the first is correct and the second is not.

\begin{theorem}[Excision is correct]\label{thm:excision}
Let $z$ be excised, giving $\Gph^{-z}$ on $V \setminus \{z\}$ with the
induced weights. Then:
\begin{enumerate}[label=(\alph*)]
\item $\Gph^{-z}$ is a contact graph in the sense of
\Cref{def:contact} --- the medium remains adjacent to every remaining
item --- so \Cref{thm:floor} holds in it with a floor $\flo^{-z}$;
\item $\sep^{-z}(v) \leq \sep(v)$ for every remaining $v$;
\item the ordering of items by separation in $\Gph^{-z}$ is not
determined by their ordering in $\Gph$, so excision is not a monotone
rescaling and must be performed before, not after, any comparison.
\end{enumerate}
\end{theorem}

\begin{proof}
(a) Deleting a non-medium vertex removes no medium edge among the
survivors, so every remaining item retains its medium edge and
\Cref{def:contact} is satisfied; \Cref{thm:floor} then applies verbatim,
with $\flo^{-z} = \min_{u \neq z} \wt(\{u,\med\}) \geq \flo$.

(b) Every admissible set for $v$ in $\Gph^{-z}$ is $S \setminus \{z\}$
for some admissible $S$ in $\Gph$, and its cut value is $\cut(S)$ minus
the weight of the $z$-edges that crossed $S$, hence no larger. Minimising
over a set of no-larger values gives $\sep^{-z}(v) \leq \sep(v)$.

(c) It suffices to exhibit two items whose order reverses. Take $v_1$
adjacent to $z$ with weight $\lambda$ and to nothing else, with
$\wt(\{v_1,\med\}) = 1$; and $v_2$ adjacent to $z$ with weight $0^{+}$
and to a private neighbour with weight $\mu$, with
$\wt(\{v_2,\med\}) = 1$. In $\Gph$, $\sep(v_1) = 1+\lambda$ and
$\sep(v_2) \approx 1 + \min(\mu,\ \cdot)$; choosing
$\lambda > \mu$ gives $\sep(v_1) > \sep(v_2)$. In $\Gph^{-z}$,
$\sep^{-z}(v_1) = 1$ while $\sep^{-z}(v_2)$ still involves $\mu$, giving
$\sep^{-z}(v_1) < \sep^{-z}(v_2)$ whenever $\mu > 0$. The order reverses.
\end{proof}

\begin{proposition}[Reweighting does not fix it]\label{prop:reweight-fails}
Let $z$'s edges be multiplied by $\alpha \in (0,1)$. Then for every
$\alpha > 0$ the conclusion of \Cref{thm:collapse} holds with $\lambda$
replaced by $\alpha\lambda$; in particular, for any $\alpha>0$ there is
a depth $d(\alpha)$ beyond which the collapse recurs. Reweighting
postpones the collapse to deeper separations rather than removing it.
\end{proposition}

\begin{proof}
The hypothesis of \Cref{thm:collapse} requires only
$\wt(\{u,z\}) \geq \lambda$ for some $\lambda>0$; scaling by $\alpha$
preserves this with $\alpha\lambda>0$. The theorem's second display then
gives $\sep(v) \geq \alpha\lambda\,\dep(v)$, and the collapse consequence
holds once $\alpha\lambda d$ dominates the subordinate terms, that is,
for $d \geq d(\alpha)$ with $d(\alpha)$ scaling as $1/\alpha$.
\end{proof}

\begin{remark}[Reading the remedy]\label{rem:remedy}
\Cref{prop:reweight-fails} says the difference between a sink and a
normal vertex is not one of degree but of kind: a sink participates in
every separation, and any positive weight on a universally-present
participant eventually dominates. \Cref{thm:excision}(c) adds the
operational sting --- because excision reorders items, it cannot be
applied as a correction to results already computed. It must be done to
the graph before any separation is reported.
\end{remark}

% =====================================================================
\section{Protocol}
\label{sec:protocol}

The results specify a four-step diagnostic, performed on every contact
graph before any separation is reported.

\begin{enumerate}[label=(\roman*)]
\item \textbf{Compute all separations and minimisers.} $n$ max-flows,
\Cref{prop:computable}. This is required anyway for the intended
analysis; the diagnostic adds nothing to it.

\item \textbf{Accumulate weighted spread in one pass.}
\Cref{prop:spread-cost}: $O(n|E|)$, one accumulator per vertex. Do
\emph{not} substitute degree --- by \Cref{thm:degree-fails} it is neither
sound nor complete.

\item \textbf{Compare against the forced threshold.} Compute $\flo$ and
$W$ from the same pass and flag every $z$ with
$\spr_{\wt}(z) > 1 - \flo/W$ (\Cref{thm:threshold}). The threshold is
read from the graph, not chosen.

\item \textbf{Excise and recompute; do not reweight, and do not correct
downstream.} By \Cref{thm:excision} excision yields a valid contact
graph with a floor; by \Cref{prop:reweight-fails} reweighting only
postpones the collapse; and by \Cref{thm:excision}(c) the item ordering
changes, so results computed before excision cannot be adjusted after it.
\end{enumerate}

\begin{remark}[Report the excision]\label{rem:report}
Which vertices were excised is part of the result, not part of the
preprocessing. Two analyses of the same data that excised different sinks
are analyses of different graphs, and their separations are not
comparable; recording the excision set is what makes the comparison
decidable.
\end{remark}

% =====================================================================
\section{Discussion}
\label{sec:discussion}

\subsection{Relation to contaminant filtering}

Standard practice filters known contaminants using curated lists
\citep{cox2008maxquant}, which is effective for the contaminants on the
list. \Cref{thm:collapse} is agnostic about the chemistry: any vertex
with universal attachment collapses the graph, whether or not it is a
known contaminant, and \Cref{def:wspread} detects it from the graph
structure alone. The two approaches are complementary --- a curated list
catches known chemistry that may have low spread; spread catches
structural pathology that may have unknown chemistry --- and neither
subsumes the other.

\subsection{Relation to hub removal in network analysis}

Removing high-degree vertices is a familiar move in network analysis.
\Cref{thm:degree-fails} is a caution against importing it directly: in a
weighted contact graph the relevant quantity is participation in minimum
cuts, not adjacency count, and the two come apart in both directions.
Where the analysis of interest is itself cut-based --- as separation is
--- the detector should be cut-based too.

\subsection{Limitations}

\Cref{thm:threshold} bounds the admissible threshold above but does not
determine it uniquely; a graph may have several vertices below the bound
whose joint excision matters, and the theorem is silent on interactions
between sinks. \Cref{prop:spread-cost} presumes the $n$ minimisers are
computed, which is the dominant cost; a sampled variant using $O(\log n)$
randomly chosen $v$ gives a concentration bound rather than an exact
spread and is not analysed here. \Cref{thm:excision}(c) is a genuine
obstruction and not an artefact: no post-hoc correction exists.

\subsection{Conclusion}

A single vertex adjacent to everything makes every separation in a
contact graph a report about that vertex rather than about the item
queried, and does so without any symptom that ordinary quality control
detects --- the outputs stay positive, stable, and reproducible.
Reproducibility is therefore not evidence of soundness in this setting.
Degree does not find such vertices. Weighted spread does, in one pass
over work already performed, against a threshold the graph itself
supplies; and the only correct remedy is excision, performed before any
separation is reported and recorded as part of the result.

% =====================================================================
\bibliographystyle{plainnat}
\bibliography{references}

\end{document}
